% O mapa da decisão: cada pregão do traço posicionado pela previsão do ARIMA
% (direção, eixo horizontal) e pela volatilidade prevista do GARCH dividida pelo
% limiar (risco, eixo vertical). A cor é a decisão que o robô de fato registrou.
\documentclass[border=6pt]{standalone}
% Preâmbulo comum dos gráficos que explicam os métodos (30/10/2026).
% Dados e números: dados/resultados/metodos_ilustrados_2026-10-30/, gerados por
% ferramentas/MetodosIlustrados.mq5 com as funções do próprio robô. Nenhum
% número é digitado aqui: os do texto vêm de valores.tex (\vX para o texto,
% \pX para coordenada).
\usepackage[T1]{fontenc}
\usepackage{amsmath}
\usepackage{mathptmx}  % Times, como as figuras do TG
\usepackage{pgfplots}
\usepgfplotslibrary{groupplots,dateplot,colormaps,fillbetween}
\usetikzlibrary{arrows.meta,positioning}
\pgfplotsset{compat=1.18}
\definecolor{cmod}{RGB}{0,90,181}     % modelo / ajuste
\definecolor{cdes}{RGB}{220,50,32}    % destaque, limite, valor crítico
\definecolor{cdad}{RGB}{120,120,120}  % dado observado
\newcommand{\arq}[1]{../../dados/resultados/metodos_ilustrados_2026-10-30/#1}
\input{\arq{valores.tex}}
% Figuras 08 a 14 (o robô operando): traço por barra da refação de 05/10/2026,
% resumido por ferramentas/scripts/robo_em_acao.sh. Macros \vR e \pR.
\newcommand{\arqr}[1]{../../dados/resultados/robo_em_acao_2026-10-30/#1}
\input{\arqr{valores.tex}}
\definecolor{cari}{RGB}{230,140,0}   % ARIMA puro
\definecolor{cok}{RGB}{0,140,70}     % compra liberada
\usetikzlibrary{shapes.misc}
\newcommand{\passo}[1]{\tikz[baseline=(n.base)]\node[circle,fill=black!75,text=white,inner sep=1pt,font=\footnotesize\bfseries](n){#1};}
% eixo de datas comum às figuras do traço (jul/2019 a jun/2021)
\pgfplotsset{traco/.style={date coordinates in=x, xmin=2019-07-01, xmax=2021-06-30,
    xtick={2019-07-01,2020-01-01,2020-07-01,2021-01-01,2021-06-30},
    xticklabels={jul/2019,jan/2020,jul/2020,jan/2021,jun/2021}}}
\pgfplotsset{
  base/.style={table/col sep=comma, grid=major, grid style={gray!20},
    tick label style={font=\small}, label style={font=\small},
    title style={font=\small\bfseries, align=center},
    every axis plot/.append style={line width=0.9pt},
    /pgf/number format/use comma, /pgf/number format/1000 sep={.}},
}
% faixa vertical nos pregões do traço em que a condição #1 vale, na cor #2,
% do piso #3 ao teto #4 do eixo
\newcommand{\faixa}[4]{%
  \addplot[name path=topo, draw=none, const plot, forget plot] table[x=data, y expr={(#1) ? #4 : #3}]{\arqr{dia_a_dia.csv}};
  \addplot[name path=piso, draw=none, forget plot] coordinates {(2019-07-01,#3) (2021-06-30,#3)};
  \addplot[#2, forget plot] fill between[of=topo and piso];}

\begin{document}
\begin{tikzpicture}
\begin{axis}[base, name=mapa, width=12cm, height=9cm, xmin=-5, xmax=5, ymin=0.4, ymax=15, ymode=log,
  log ticks with fixed point, ytick={0.5,1,2,4,8}, yticklabels={0{,}5,1,2,4,8},
  xlabel={\passo{A} ARIMA: retorno previsto para o dia (\%) $\;\Rightarrow\;$ direção},
  ylabel={\passo{G} GARCH: $\sigma$ prevista $\div$ limiar $\;\Rightarrow\;$ risco},
  title={Cada ponto é um pregão (\vRInicio{} a \vRFim)}, clip=true]
\fill[cdes!10] (axis cs:-5,1) rectangle (axis cs:5,15);
\fill[cok!12] (axis cs:0,0.4) rectangle (axis cs:5,1);
\fill[black!7] (axis cs:-5,0.4) rectangle (axis cs:0,1);
\draw[black!70, line width=1pt] (axis cs:-5,1) -- (axis cs:5,1);
\draw[black!70, line width=1pt] (axis cs:0,0.4) -- (axis cs:0,1);
\draw[black!30, dashed] (axis cs:0,1) -- (axis cs:0,15);
\addplot[only marks, mark=*, mark size=1.3pt, cok, opacity=0.75] table[x=prev, y=razao]{\arqr{decisao_compra.csv}};
\addplot[only marks, mark=*, mark size=1.3pt, black!55, opacity=0.75] table[x=prev, y=razao]{\arqr{decisao_venda.csv}};
\addplot[only marks, mark=*, mark size=1.3pt, cdes, opacity=0.75] table[x=prev, y=razao]{\arqr{decisao_vetada.csv}};
\node[cdes, font=\small\bfseries, align=center, fill=white, fill opacity=0.8, text opacity=1] at (axis cs:-3.2,9)
  {VETADA (\vRVetos)\\[-1pt]\normalfont\footnotesize não compra e zera a posição};
\node[cok!70!black, font=\small\bfseries, align=center, fill=white, fill opacity=0.8, text opacity=1] at (axis cs:3,0.5)
  {COMPRA (\vRCompras)\\[-1pt]\normalfont\footnotesize compra, ou mantém a compra};
\node[black!70, font=\small\bfseries, align=center, fill=white, fill opacity=0.8, text opacity=1] at (axis cs:-3,0.5)
  {VENDA (\vRVendas)\\[-1pt]\normalfont\footnotesize o robô só compra: zera a compra};
\node[font=\footnotesize, anchor=south east, black!70] at (axis cs:4.95,1) {$\sigma$ prevista $=$ limiar};
\end{axis}
\node[font=\footnotesize, text=black!75, align=left, anchor=north west, text width=11cm]
  at ([yshift=-0.2cm]mapa.outer south west) {Como ler: o ARIMA decide a direção (esquerda ou direita da linha vertical) e o
  GARCH decide se o risco permite operar (abaixo ou acima da linha horizontal). O GARCH nunca muda a direção, só veta.
  Sem o GARCH, os \vRVetoCompra{} pontos vermelhos à direita seriam compras e os \vRVetoVenda{} à esquerda seriam vendas
  (o ARIMA puro teve \vRComprasA{} sinais de compra e \vRVendasA{} de venda). Conferência: a regra
  ``veta se $\sigma$ prevista $>$ limiar'', refeita a partir das colunas do traço, coincide com o veto registrado pelo robô
  nos \vRN{} pregões (\vRIncoerentes{} diferenças).};
\end{tikzpicture}
\end{document}
